\section{Preliminary Simulation Results}
\label{sec:prelim}

This chapter presents the simulation environment configured
for evaluating the DMAP-SDVN framework, the implementation
status of the two completed tasks, and evidence from
preliminary simulation runs conducted to validate the
attack injection and metric collection infrastructure
prior to the implementation of the full signature
detection layer in Task~03.

% ─────────────────────────────────────────────────────────
\subsection{Simulation Environment and Configuration}

% ─────────────────────────────────────────────────────────

The evaluation environment is built on NS3 version~3.35
with the OFSwitch13 module providing OpenFlow~1.3
control-plane support. The simulation models an SDVN
topology with a logically centralised SDN controller,
multiple RSU nodes, and mobile vehicular nodes
communicating over CSMA downlink and LTE uplink
channels, reflecting the V2I and V2V communication
paths described in Section~\ref{sec:system_model}.

Table~\ref{tab:sim_params} summarises the primary
simulation parameters used in the current runs. The
\texttt{attack\_percentage} parameter controls the
fraction of total network nodes designated as
attackers; the current preliminary runs use 40\%
attack intensity, yielding seven attacking nodes.

\begin{table}[htb]
\centering
\caption{Simulation Parameters (Current Implementation)}
\label{tab:sim_params}
\setlength{\tabcolsep}{5pt}
\begin{tabular}{ll}
\hline
\textbf{Parameter} & \textbf{Value} \\
\hline
NS3 version                      & 3.35 \\
Control-plane protocol           & OpenFlow~1.3 (OFSwitch13) \\
Channel types                    & CSMA (downlink), LTE (uplink) \\
Simulation duration              & 10\,s per run \\
Metrics epoch interval           & 1\,s (periodic CSV flush) \\
Attack variants implemented      & DI, RC, SC, SE \\
Attack intensity (current runs)  & 40\% of total nodes \\
Number of attacking nodes        & 7 \\
DI injected delays               & 50\,ms ($t{=}1$\,s), 80\,ms ($t{=}4$\,s),
                                   120\,ms ($t{=}7$\,s) \\
RC trigger times                 & $t{=}2$\,s, 5\,s, 8\,s \\
SC observation times             & $t{=}1.5$, 3.5, 5.5, 7.5, 9.5\,s \\
SE beacon gap times              & $t{=}3$\,s, 6\,s, 9\,s \\
SC leakage threshold             & 5 observations per node \\
SE oscillation threshold         & 2 missing beacons per node \\
Trust score decrement per event  & 0.03 \\
Initial RSU trust score          & 1.0 (fully trusted) \\
Trust quarantine threshold       & 0.4 \\
Performance metrics logged       & ADR (per variant), FPR, MCC, TTM, TSA \\
Output files                     & CSV (per-epoch and per-event) \\
\hline
\end{tabular}
\end{table}

As shown in Table~\ref{tab:sim_params}, the simulation
is configured with four attack variants injected at
40\% node intensity across a 10\,s simulation window,
with all performance metrics flushed to CSV every
1\,s epoch.

% ─────────────────────────────────────────────────────────
\subsection{Task~01: Performance Metric Infrastructure}
% ─────────────────────────────────────────────────────────

Task~01 (completed) implemented the full performance
metric collection infrastructure within the NS3
simulation. All six metrics defined in
Section~\ref{sec:baselines} are computed and written
to CSV at every 1\,s epoch interval by the function
\texttt{dmap\_periodic\_flush}, which is scheduled
using \texttt{Simulator::Schedule} at simulation
start.

The metric state maintained globally throughout the
simulation is as follows. Per-variant Attack Detection
Rate counters (\texttt{detected\_DI},
\texttt{detected\_RC}, \texttt{detected\_SC},
\texttt{detected\_SE}) track correctly detected
attacks against the total injected count per variant.
False Positive Rate accumulators track benign events
incorrectly flagged. Time-to-Mitigate records the
elapsed time between attack confirmation and
enforcement of the quarantine or rerouting action.
Trust Score Accuracy is derived from a per-RSU
confusion matrix classifying each RSU as correctly
high-trust, correctly low-trust, false-positive
low-trust, or missed-compromise. Matthews Correlation
Coefficient is computed from the aggregate confusion
matrix across all variants. Per-RSU trust scores are
initialised to~1.0 and decremented by~0.03 per
attack event, flooring at~0.

\subsubsection{Output File: \texttt{dmap\_epoch\_metrics.csv}}

This file contains one row per simulation epoch and
records the complete metric state at that moment.
Table~\ref{tab:epoch_schema} describes all columns.

\begin{table}[htb]
\centering
\caption{Column Schema of \texttt{dmap\_epoch\_metrics.csv}}
\label{tab:epoch_schema}
\scriptsize
\setlength{\tabcolsep}{3pt}
\begin{tabular}{lll}
\hline
\textbf{Column} & \textbf{Type} & \textbf{Description} \\
\hline
\texttt{sim\_time\_sec}              & float  &
  Simulation clock time in seconds \\
\texttt{epoch\_id}                   & int    &
  Epoch index (0-based) \\
\texttt{mode}                        & string &
  Defence mode active: LW or FULL \\
\texttt{MCC\_DI/RC/SC/SE/ALL}        & float  &
  Matthews Correlation Coefficient per variant and overall \\
\texttt{ADR\_DI/RC/SC/SE/ALL}        & float  &
  Attack Detection Rate per variant and overall \\
\texttt{FPR\_DI/RC/SC/SE/ALL}        & float  &
  False Positive Rate per variant and overall \\
\texttt{latency\_overhead\_LW\_ms}   & float  &
  Control-plane latency overhead in LW mode (ms) \\
\texttt{latency\_overhead\_FULL\_ms} & float  &
  Control-plane latency overhead in FULL mode (ms) \\
\texttt{TTM\_DI/RC/SC/SE/ALL\_ms}    & float  &
  Average Time-to-Mitigate per variant (ms) \\
\texttt{TSA\_precision}              & float  &
  Fraction of low-trust RSUs that were truly compromised \\
\texttt{TSA\_recall}                 & float  &
  Fraction of compromised RSUs that received low trust \\
\texttt{TSA\_F1}                     & float  &
  Harmonic mean of TSA precision and recall \\
\texttt{TP\_ALL / TN\_ALL}           & int    &
  True positive and true negative counts (all variants) \\
\texttt{FP\_ALL / FN\_ALL}           & int    &
  False positive and false negative counts (all variants) \\
\hline
\end{tabular}
\end{table}

As shown in Table~\ref{tab:epoch_schema}, the epoch
metrics file captures 30 columns per row, covering
per-variant and aggregate values for all five
performance metrics together with the raw confusion
matrix counts needed to reproduce any derived metric.

\subsubsection{Output File: \texttt{dmap\_degradation.csv}}

This file records a condensed metric snapshot at each
epoch alongside the configured attack intensity
percentage. Table~\ref{tab:degradation_schema}
describes its columns.

\begin{table}[htb]
\centering
\caption{Column Schema of \texttt{dmap\_degradation.csv}}
\label{tab:degradation_schema}
\setlength{\tabcolsep}{4pt}
\begin{tabular}{lll}
\hline
\textbf{Column} & \textbf{Type} & \textbf{Description} \\
\hline
\texttt{sim\_time\_ms} & int   &
  Epoch simulation time in milliseconds \\
\texttt{attack\_pct}   & int   &
  Configured attack intensity percentage \\
\texttt{ADR\_DI}       & float &
  Attack Detection Rate for Delay Injection \\
\texttt{ADR\_RC}       & float &
  Attack Detection Rate for Race Condition \\
\texttt{ADR\_SC}       & float &
  Attack Detection Rate for Timing Side-Channel \\
\texttt{ADR\_SE}       & float &
  Attack Detection Rate for Synchronisation Exploit \\
\texttt{FPR\_ALL}      & float &
  Aggregate False Positive Rate \\
\texttt{MCC\_ALL}      & float &
  Aggregate Matthews Correlation Coefficient \\
\texttt{avg\_TTM\_ms}  & float &
  Average Time-to-Mitigate across all variants (ms) \\
\hline
\end{tabular}
\end{table}

As shown in Table~\ref{tab:degradation_schema}, this
file stores nine columns per row and is designed to
support cross-run comparison of metric quality as
attack intensity is swept from 0 to 100\% across
multiple simulation runs, enabling the degradation
curves to be generated in the final evaluation chapter.

% ─────────────────────────────────────────────────────────
\subsection{Task~02: Attack Scenario Implementation}
% ─────────────────────────────────────────────────────────

Task~02 (completed) implemented all four timing attack
injection functions within \texttt{LDA.cc} and
configured a degradation snapshot recorder. The four
functions 
\texttt{inject\_de\-lay\_in\-jec\-tion\_at\-tack},
\texttt{inject\_race\_con\-di\-tion\_at\-tack},
\texttt{inject\_side\_chan\-nel\_at\-tack}, and
\texttt{inject\_sync\_ex\-ploit\_at\-tack}
map to the existing malicious-node flag arrays
(\texttt{flood\-ing\_ma\-li\-cious\_nodes} for DI,
\texttt{fab\-ri\-ca\-tion\_ma\-li\-cious\_nodes} for RC,
\texttt{MIM\_ma\-li\-cious\_nodes} for SC, and
\texttt{van\-ish\-ing\_ma\-li\-cious\_nodes} for SE),
ensuring compatibility with the existing routing and
detection infrastructure.

The number of attacking nodes is computed as:
\begin{equation*}
N_{\mathrm{attack}} =
\left\lceil
  \frac{\mathtt{attack\_percentage}}{100}
  \times \mathtt{total\_size}
\right\rceil
\end{equation*}

Table~\ref{tab:attack_schedule} shows the full
injection schedule per attacking node.

\begin{table}[htb]
\centering
\caption{Attack Injection Schedule per Attacking Node
(Task~02, 40\% intensity)}
\label{tab:attack_schedule}
\setlength{\tabcolsep}{4pt}
\begin{tabular}{llll}
\hline
\textbf{Attack} & \textbf{Trigger Times (s)} &
\textbf{Parameter logged} & \textbf{Effect modelled} \\
\hline
DI & 1.0, 4.0, 7.0 &
  Injected delay (s) & Flow-Mod message buffered \\
RC & 2.0, 5.0, 8.0 &
  0 (flag only) & Contradictory link events \\
SC & 1.5, 3.5, 5.5, 7.5, 9.5 &
  Observation count & Passive timing corpus grows \\
SE & 3.0, 6.0, 9.0 &
  Beacon gap count & Handover oscillation triggered \\
\hline
\end{tabular}
\end{table}

As shown in Table~\ref{tab:attack_schedule}, the four
attack variants are staggered across the 10\,s window
to produce overlapping but distinct attack events,
reflecting a realistic scenario where multiple attack
types co-exist in the same network simultaneously.

\subsubsection{Console Log Evidence}

Figure~\ref{fig:log1} presents the NS3 simulation
console output at $t{=}1$\,s and $t{=}2$\,s.

\begin{figure}[htb]
\centering
\includegraphics[width=\linewidth]
{Progress_Works/NS3-Images/sim_log1.png}
\caption{NS3 simulation console output at $t{=}1$\,s
and $t{=}2$\,s.}
\label{fig:log1}
\end{figure}
\FloatBarrier

As shown in Figure~\ref{fig:log1}, normal CSMA
downlink and LTE uplink packet events precede the
first attack injection. At $t{=}1$\,s,
\texttt{inject\_delay\_injection\_attack} fires for
all seven attacking nodes with 50\,ms injected delay,
setting \texttt{flooding\_malicious\_nodes[i] = true}
and decrementing each node's trust score by~0.03.
At $t{=}1.5$\,s the first Side-Channel observation
is logged for each node, incrementing
\texttt{sc\_observation\_count} toward the leakage
threshold of~5. At $t{=}2$\,s Race Condition events
are triggered across all attacking nodes, setting
\texttt{rc\_race\_active[i] = true} for each node.

Figure~\ref{fig:log2} presents the NS3 simulation
console output at $t{=}3$\,s to $t{=}4$\,s.

\begin{figure}[htb]
\centering
\includegraphics[width=\linewidth]
{Progress_Works/NS3-Images/sim_log2.png}
\caption{NS3 simulation console output at $t{=}3$\,s
to $t{=}4$\,s.}
\label{fig:log2}
\end{figure}
\FloatBarrier

As shown in Figure~\ref{fig:log2}, at $t{=}3$\,s
the first Synchronisation Exploit fires for each
attacking node, incrementing
\texttt{se\_beacon\_gap[i]} toward the oscillation
threshold of~2. SC observation~\#2 follows at
$t{=}3.5$\,s. At $t{=}4$\,s the second DI wave
increases the injected delay to 80\,ms, modelling
an escalating attacker who progressively extends the
buffering window to produce larger temporal deviations.

Figure~\ref{fig:log3} presents the NS3 simulation
console output at $t{=}5$\,s to $t{=}5.5$\,s.

\begin{figure}[htb]
\centering
\includegraphics[width=\linewidth]
{Progress_Works/NS3-Images/sim_log3.png}
\caption{NS3 simulation console output at $t{=}5$\,s
to $t{=}5.5$\,s.}
\label{fig:log3}
\end{figure}
\FloatBarrier

As shown in Figure~\ref{fig:log3}, Race Condition
events fire a second time at $t{=}5$\,s across all
attacking nodes. SC observation~\#3 is logged at
$t{=}5.5$\,s, continuing the accumulation of the
passive timing corpus. The progressive growth of the
observation count across five trigger points reflects
the enlarging statistical corpus that enables the
attacker to infer control-plane timing patterns,
corresponding to the side-channel amplification factor
$\Phi_{\mathrm{SC}}(t)$ described in
Section~\ref{sec:mobility_amp}.

\subsubsection{Output File: \texttt{dmap\_attack\_scenario\_log.csv}}

This file records every individual attack injection
event produced during the simulation.
Table~\ref{tab:attack_log_schema} describes its
columns.

\begin{table}[htb]
\centering
\caption{Column Schema of
\texttt{dmap\_attack\_scenario\_log.csv}}
\label{tab:attack_log_schema}
\setlength{\tabcolsep}{4pt}
\begin{tabular}{lll}
\hline
\textbf{Column} & \textbf{Type} & \textbf{Description} \\
\hline
\texttt{time\_ms}     & int    &
  Simulation time of injection event (ms) \\
\texttt{node\_id}     & int    &
  Index of the attacked RSU node \\
\texttt{attack\_type} & string &
  Attack variant: DI, RC, SC, or SE \\
\texttt{param}        & float  &
  Injected delay in seconds (DI); observation
  count (SC); beacon gap count (SE); 0 for RC \\
\texttt{attack\_pct}  & int    &
  Configured attack intensity percentage \\
\hline
\end{tabular}
\end{table}

Table~\ref{tab:attack_log_sample} shows a
representative sample of rows from the 40\% attack
run.

\begin{table}[htb]
\centering
\caption{Sample Rows from
\texttt{dmap\_attack\_scenario\_log.csv}
(40\% attack intensity)}
\label{tab:attack_log_sample}
\scriptsize
\setlength{\tabcolsep}{5pt}
\begin{tabular}{rrllr}
\hline
\textbf{time\_ms} & \textbf{node\_id} &
\textbf{attack\_type} & \textbf{param} &
\textbf{attack\_pct} \\
\hline
1000 & 0 & DI & 0.050 & 40 \\
1000 & 1 & DI & 0.050 & 40 \\
1000 & 2 & DI & 0.050 & 40 \\
1000 & 3 & DI & 0.050 & 40 \\
1000 & 4 & DI & 0.050 & 40 \\
1000 & 5 & DI & 0.050 & 40 \\
1000 & 6 & DI & 0.050 & 40 \\
1500 & 0 & SC & 1     & 40 \\
1500 & 1 & SC & 1     & 40 \\
2000 & 0 & RC & 0     & 40 \\
3000 & 0 & SE & 1     & 40 \\
3500 & 0 & SC & 2     & 40 \\
4000 & 0 & DI & 0.080 & 40 \\
5500 & 0 & SC & 3     & 40 \\
6000 & 0 & SE & 2     & 40 \\
7000 & 0 & DI & 0.120 & 40 \\
7500 & 0 & SC & 4     & 40 \\
9000 & 0 & SE & 3     & 40 \\
9500 & 0 & SC & 5     & 40 \\
\hline
\end{tabular}
\end{table}

As shown in Table~\ref{tab:attack_log_sample}, all
seven attacking nodes receive DI events at
$t{=}1000$\,ms with 50\,ms injected delay. SC
observations accumulate from observation~\#1 at
$t{=}1500$\,ms through observation~\#5 at
$t{=}9500$\,ms, reaching the leakage threshold of~5
by end of simulation. The DI injected delay escalates
from 50\,ms to 80\,ms to 120\,ms across the three
trigger points, and SE beacon gaps reach the
oscillation threshold of~2 by $t{=}6000$\,ms.

% ─────────────────────────────────────────────────────────
\subsection{Current Metric Values and Interpretation}
% ─────────────────────────────────────────────────────────

Inspection of both \texttt{dmap\_degradation.csv} and
\texttt{dmap\_epoch\_metrics.csv} from the current
preliminary runs shows that all metric values ---
ADR, FPR, MCC, TTM, and TSA --- remain at~0.0 across
all epochs at the current implementation stage. This
is the expected and correct result, and it reflects
an important architectural separation within the
DMAP-SDVN implementation.

The metric values are zero because the signature
detection layer has not yet been connected to the
confusion matrix update logic. Task~02 injects the
four attack types and logs every event, but the
classification of each event as a True Positive,
False Negative, True Negative, or False Positive
requires the four signature indicators
$S_{\mathrm{DI}}$, $S_{\mathrm{RC}}$,
$S_{\mathrm{SC}}$, $S_{\mathrm{SE}}$ from
Section~\ref{sec:signatures} to be evaluated and
used to increment the TP, TN, FP, FN counters inside
\texttt{dmap\_periodic\_flush}. This connection is
the primary objective of Task~03.

The zero-metric result is itself meaningful: it
quantifies the performance of the \emph{undefended}
system --- a network that receives all four attack
types but has no active detection mechanism. This
constitutes the pre-DMAP baseline against which the
defended framework will be compared in the final
evaluation. Table~\ref{tab:zero_baseline} summarises
this baseline state.

\begin{table}[htb]
\centering
\caption{Undefended Baseline Metric Values
(Task~02 Only, No Signature Detection Active)}
\label{tab:zero_baseline}
\setlength{\tabcolsep}{4pt}
\begin{tabular}{lrl}
\hline
\textbf{Metric} & \textbf{Value} &
\textbf{Interpretation} \\
\hline
ADR (all variants) & 0.00 &
  No attacks detected --- zero true positives \\
FPR                & 0.00 &
  No benign events flagged --- zero false positives \\
MCC                & 0.00 &
  No classification decisions made \\
avg TTM            & 0\,ms &
  No mitigation triggered \\
TSA Precision      & 0.00 &
  No trust-based isolation decisions \\
TSA Recall         & 0.00 &
  All compromised nodes missed by detection \\
TP / TN / FP / FN  & 0 / 0 / 0 / 0 &
  Confusion matrix empty pending Task~03 \\
\hline
\end{tabular}
\end{table}

As shown in Table~\ref{tab:zero_baseline}, the
complete absence of detection activity confirms that
the simulation is operating as an undefended baseline.
This result will serve as the reference point for the
Task~03 evaluation, where the DMAP-SDVN signature
detection layer is expected to produce non-zero ADR,
MCC, and TTM values demonstrating measurable detection
and mitigation capability.

% ─────────────────────────────────────────────────────────
\subsection{Trust Score Degradation Under Attack}
% ─────────────────────────────────────────────────────────

Although detection metrics remain at zero pending
Task~03, the trust score degradation produced by the
attack injection functions is active and observable.
Each attack event decrements the attacked node's
trust score by~0.03 as implemented in the injection
functions and consistent with the trust update rule
in Section~\ref{sec:blockchain}.

Across the 10\,s simulation at 40\% intensity, each
attacking node receives 3~DI events, 3~RC events,
5~SC observations, and 3~SE events, totalling
14~trust decrements of~0.03 each. This reduces each
node's trust score from $\psi_r = 1.0$ to:
\begin{equation*}
\psi_r = 1.0 - (14 \times 0.03) = 0.58
\end{equation*}

This value remains above the quarantine threshold
$\psi_{\mathrm{low}} = 0.4$, confirming that the
current 40\% attack intensity produces measurable
trust degradation without yet triggering automatic
RSU isolation. At higher attack intensities to be
evaluated in Task~03, more frequent events will push
trust scores below $\psi_{\mathrm{low}}$, activating
the smart-contract quarantine path and producing
non-zero TTM values.

Figure~\ref{fig:attack_timeline} presents the
complete attack event timeline and the corresponding
per-node trust score degradation trajectory over the
10\,s simulation window.

\begin{figure}[htb]
\centering
\includegraphics[width=\linewidth]
{Progress_Works/NS3-Images/dmap_attack_timeline.png}
\caption{Attack event timeline and trust score
degradation across seven attacking nodes under
40\% attack intensity over a 10\,s simulation window.}
\label{fig:attack_timeline}
\end{figure}
\FloatBarrier

As shown in Figure~\ref{fig:attack_timeline}, the
upper panel confirms that all four attack variants
--- DI, RC, SC, and SE --- are injected at the
scheduled times across all seven nodes, with DI
delay escalating from 50\,ms at $t{=}1$\,s to
80\,ms at $t{=}4$\,s and 120\,ms at $t{=}7$\,s.
The lower panel shows the step-wise trust score
trajectory for a representative node, decreasing
monotonically from~1.0 to~0.58 across 14 events,
remaining above $\psi_{\mathrm{low}} = 0.40$
throughout the simulation and thereby confirming
that no automatic quarantine is triggered at
40\% attack intensity.

% ─────────────────────────────────────────────────────────
\subsection{Implementation Status and Next Steps}
% ─────────────────────────────────────────────────────────

Table~\ref{tab:task_status} summarises the current
implementation status against the four tasks assigned
by the supervisor.

\begin{table}[htb]
\centering
\caption{Implementation Task Status}
\label{tab:task_status}
\begin{tabularx}{\linewidth}{l X l} % Adjusted to 3 columns: Task (l), Description (X), Status (l)
\hline
\textbf{Task} & \textbf{Description} & \textbf{Status} \\
\hline
01 & Implement all PEMs and save to CSV periodically & \textbf{Complete} \\
02 & Implement all four attack scenarios and establish undefended baseline & \textbf{Complete} \\
03 & Implement full signature detection for all four variants and show metric improvement & In progress \\
04 & Complete result figures and finalise report for supervisor signature & Pending \\
\hline
\end{tabularx}
\end{table}

As shown in Table~\ref{tab:task_status}, Tasks~01
and~02 are complete. Task~03 will implement the four
attack signature indicators from
Section~\ref{sec:signatures} inside the NS3
simulation and connect them to the confusion matrix
update logic within \texttt{dmap\_periodic\_flush}.
Once connected, the TP, TN, FP, and FN counters will
accumulate correctly per epoch, producing non-zero
ADR, MCC, FPR, and TTM values. Comparison of these
defended metric values against the zero-valued
undefended baseline from Task~02 will directly
quantify the detection improvement delivered by the
DMAP-SDVN signature detection layer. Task~04 will
then complete all remaining result figures and
submit the final report for supervisor review and
signature.